% General artificial intelligence terms, in alphabetical order.
\domain{Artificial Intelligence}

\begin{entry}{agent}{Agent}
  \defn*{russell-norvig-2021}{An agent is anything that can be viewed
    as perceiving its environment through sensors and acting upon that
    environment through actuators.}
  \defn{wooldridge-1995}{A hardware or, more usually, software-based computer
    system with four properties: autonomy (it operates without direct human
    intervention), social ability (it interacts with other agents or humans),
    reactivity (it perceives and responds to changes in its environment), and
    pro-activeness (it takes the initiative to pursue goals).}
  \defn{iso-22989}{An automated entity that senses and responds to its
    environment and takes actions to achieve its goals.}
  \note{Russell and Norvig's definition is deliberately broad and includes a
    thermostat. Wooldridge and Jennings add requirements, such as autonomy and
    goal-directedness, that separate agents from ordinary programs.}
  \related{agentic-ai,autonomy,ai-system}
\end{entry}

\begin{entry}{agentic-ai}{Agentic AI}[AI agent]
  \defn{shavit-2023}{AI systems that can pursue complex goals with limited
    direct supervision, including by taking actions in the world that were
    not specified step by step by a human.}
  \defn{chan-2023}{Algorithmic systems that show increasing degrees of
    agency. Agency has four characteristics: goals are given without
    specifying how to reach them; actions have an impact without a human
    intermediary; behaviour is goal-directed; and the system plans over the
    long term.}
  \defn{owasp-llm-2025}{LLM-based systems granted \emph{agency}: the ability
    to call functions, use tools, or interface with other systems in
    response to prompts. \enquote{Excessive agency} is the risk that such
    functionality, permissions, or autonomy go beyond what the task needs.}
  \note{\enquote{Agentic} is best understood as a matter of degree rather
    than a binary property. Governance frameworks usually focus on the
    delegation of actions, and on who is accountable for them.}
  \related{agent,autonomy,human-oversight,prompt-injection}
\end{entry}

\begin{entry}{agi}{Artificial General Intelligence}[AGI]
  \defn{goertzel-2007}{AI systems with the ability to solve a wide variety of
    problems in a wide variety of contexts, including problems they were not
    specifically designed for, as opposed to systems built for one narrow
    task.}
  \defn{morris-2024}{A system whose capabilities match or exceed those of
    skilled humans across a wide range of non-physical tasks, including the
    ability to learn new tasks. AGI is best characterised by levels of
    performance and generality, not a single threshold.}
  \defn{searle-1980}{Related to \enquote{strong AI}: the claim that an
    appropriately programmed computer does not merely simulate a mind but
    literally has cognitive states and understanding.}
  \note{No agreed operational definition exists. Searle's strong AI is a
    philosophical claim about minds, whereas modern usage, as in Morris et al.,
    is about measurable capability.}
  \related{narrow-ai,artificial-intelligence}
\end{entry}

\begin{entry}{artificial-intelligence}{Artificial Intelligence}[AI]
  \defn*{mccarthy-2007}{It is the science and engineering of making
    intelligent machines, especially intelligent computer programs.}
  \defn*{nilsson-2010}{Artificial intelligence is that activity
    devoted to making machines intelligent, and intelligence is that quality
    that enables an entity to function appropriately and with foresight in
    its environment.}
  \defn{russell-norvig-2021}{The study and construction of agents that act
    rationally, meaning that they do the right thing given what they know,
    as opposed to the study of systems that think or act like humans.}
  \defn{iso-22989}{The research and development of mechanisms and
    applications of AI systems. The term also names the discipline itself.}
  \note{The earlier academic definitions describe a \emph{field of study},
    while standards and regulation mostly define the \emph{AI system}. When
    applicability or compliance is at stake, use the system-level definitions.}
  \related{ai-system,machine-learning,agent}
\end{entry}

\begin{entry}{ai-system}{AI System}
  \defn*[Art.~3(1)]{eu-ai-act}{A machine-based system that is designed to
    operate with varying levels of autonomy and that may exhibit
    adaptiveness after deployment, and that, for explicit or implicit
    objectives, infers, from the input it receives, how to generate outputs
    such as predictions, content, recommendations, or decisions that can
    influence physical or virtual environments.}
  \defn*{oecd-2024}{A machine-based system that, for explicit or implicit
    objectives, infers, from the input it receives, how to generate outputs
    such as predictions, content, recommendations, or decisions that can
    influence physical or virtual environments. Different AI systems vary in
    their levels of autonomy and adaptiveness after deployment.}
  \defn*{nist-ai-rmf}{An engineered or machine-based system that can,
    for a given set of objectives, generate outputs such as predictions,
    recommendations, or decisions influencing real or virtual environments.}
  \defn{iso-22989}{An engineered system that generates outputs such as
    content, forecasts, recommendations, or decisions for a given set of
    human-defined objectives.}
  \note{The EU and OECD definitions were aligned in 2023–2024 and both make
    \emph{inference} the distinguishing feature. That excludes conventional
    software whose rules are defined only by humans. The NIST and ISO/IEC
    definitions predate this change and do not use the inference test.}
  \related{artificial-intelligence,autonomy,general-purpose-ai}
\end{entry}

\begin{entry}{autonomy}{Autonomy}
  \defn[Recital~12]{eu-ai-act}{The degree to which an AI system's actions
    are independent of human involvement, and its ability to operate without
    human intervention.}
  \defn{iso-22989}{A characteristic of a system that can modify its intended
    domain of use or goals without external intervention, control, or
    oversight. Levels of automation run from none, through heteronomous
    systems under human control, to autonomous systems.}
  \defn{wooldridge-1995}{The ability of an agent to operate without the
    direct intervention of humans or others, having some control over its
    own actions and internal state.}
  \note{ISO/IEC 22989 separates \emph{automation}, which is carrying out a
    human-defined task, from \emph{autonomy}, which includes changing goals.
    Much everyday use of \enquote{autonomous} refers only to high automation.}
  \related{agent,agentic-ai,human-oversight}
\end{entry}

\begin{entry}{foundation-model}{Foundation Model}
  \defn*{bommasani-2021}{Any model that is trained on broad data
    (generally using self-supervision at scale) that can be adapted (e.g.,
    fine-tuned) to a wide range of downstream tasks.}
  \defn{nist-ai-600-1}{Discussed via the \enquote{dual-use foundation
    model}: an AI model trained on broad data, generally using
    self-supervision, with at least tens of billions of parameters, that is
    applicable across a wide range of contexts and could pose serious risks
    to security, public health, or safety.}
  \defn[Art.~3(63)]{eu-ai-act}{Treated in the regulation as a
    \enquote{general-purpose AI model}. See that entry.}
  \note{The term was coined to stress that one model is a shared, and
    therefore single-point-of-failure, basis for many applications.
    Regulators generally use \enquote{general-purpose AI} instead.}
  \related{general-purpose-ai,large-language-model,fine-tuning,self-supervised-learning}
\end{entry}

\begin{entry}{general-purpose-ai}{General-Purpose AI Model}[GPAI]
  \defn[Art.~3(63)]{eu-ai-act}{An AI model, including one trained with a
    large amount of data using self-supervision at scale, that displays
    significant generality, is capable of competently performing a wide range
    of distinct tasks, and can be integrated into a variety of downstream
    systems or applications. Models used only for research, development, or
    prototyping before being placed on the market are excluded.}
  \defn[Art.~3(66)]{eu-ai-act}{A \emph{general-purpose AI system} is an AI
    system based on a general-purpose AI model that can serve a variety of
    purposes, both for direct use and for integration into other AI systems.}
  \note{The EU AI Act separates the \emph{model} from the \emph{system} built
    on it and assigns obligations to each. It adds further obligations for
    models with \enquote{systemic risk}.}
  \related{foundation-model,ai-system}
\end{entry}

\begin{entry}{generative-ai}{Generative AI}[GenAI]
  \defn{nist-ai-600-1}{The class of AI models that emulate the structure and
    characteristics of input data in order to generate derived synthetic
    content, such as images, videos, audio, text, and other digital content.}
  \defn[ch.~20]{goodfellow-2016}{Models that learn a probability
    distribution over data, explicitly or implicitly, and can draw new
    samples from it. Examples are variational autoencoders, generative
    adversarial networks, and autoregressive models.}
  \defn[Art.~50]{eu-ai-act}{AI systems, including general-purpose AI
    systems, that generate synthetic audio, image, video, or text content.
    Their providers must mark outputs as artificially generated in a
    machine-readable format.}
  \note{The policy definitions refer to outputs (synthetic content). The
    technical definition refers to what the model learns (a data
    distribution).}
  \related{large-language-model,foundation-model,hallucination}
\end{entry}

\begin{entry}{hallucination}{Hallucination}[confabulation]
  \defn{ji-2023-hallucination}{Generated content that is nonsensical or
    unfaithful to the provided source content. \emph{Intrinsic}
    hallucinations contradict the source, and \emph{extrinsic}
    hallucinations cannot be verified from it.}
  \defn{nist-ai-600-1}{Termed \emph{confabulation}: the production of
    confidently stated but erroneous or false content, which may mislead or
    deceive users.}
  \defn{zhao-2023}{Content generated by a large language model that is
    factually incorrect or inconsistent with its input or context, but
    presented fluently and plausibly.}
  \note{NIST prefers \enquote{confabulation} to avoid attributing perception
    to the system. Research usage separates unfaithfulness to a given source
    from factual error about the world.}
  \related{large-language-model,retrieval-augmented-generation,validity-reliability}
\end{entry}

\begin{entry}{large-language-model}{Large Language Model}[LLM]
  \defn{zhao-2023}{A Transformer-based language model with a very large
    number of parameters, typically billions or more, that is pre-trained on
    massive text corpora and shows abilities, such as in-context learning,
    that are not found in smaller models.}
  \defn{brown-2020}{An autoregressive language model that, at sufficient
    scale, can perform new tasks from a natural-language instruction or a few
    examples in its prompt, without updating its weights.}
  \defn{bommasani-2021}{The most prominent kind of foundation model: a
    language model trained on broad text data that is adapted, by
    fine-tuning or prompting, to many downstream language tasks.}
  \note{\enquote{Large} is relative and has shifted over time. The defining
    features in practice are pre-training at scale and general-purpose
    prompting.}
  \related{foundation-model,transformer,prompt,hallucination}
\end{entry}

\begin{entry}{narrow-ai}{Narrow AI}[weak AI]
  \defn{goertzel-2007}{AI systems designed and trained for a specific task or
    a narrow range of tasks, which cannot generalise outside that domain
    without being redesigned or retrained.}
  \defn{searle-1980}{\enquote{Weak AI}: the position that computers are
    powerful tools for studying and simulating the mind, without any claim
    that the simulation itself has understanding.}
  \defn{morris-2024}{Systems with high performance on a clearly scoped task,
    such as a chess engine or a protein-structure predictor, placed on the
    \emph{narrow} end of a generality scale.}
  \related{agi,artificial-intelligence}
\end{entry}

\begin{entry}{prompt}{Prompt}
  \defn{liu-2023-prompt}{The input text, often a template with slots to be
    filled, that is given to a pre-trained language model so that the
    desired task is framed as text the model can continue or complete.}
  \defn{nist-ai-100-2}{The input to a generative AI model. In deployed
    applications it combines a \emph{system prompt} written by the
    application developer with user input and, often, content retrieved from
    external resources.}
  \defn{brown-2020}{The context given to a model at inference time, which may
    contain a task description and zero, one, or a few demonstrations, as in
    \emph{zero-shot}, \emph{one-shot}, and \emph{few-shot} prompting.}
  \related{large-language-model,prompt-injection}
\end{entry}

\begin{entry}{retrieval-augmented-generation}{Retrieval-Augmented Generation}[RAG]
  \defn{lewis-2020}{A method that combines a pre-trained parametric
    generator with a non-parametric memory, such as a dense vector index of
    documents. Retrieved passages condition the generation.}
  \defn{gao-2023-rag}{A paradigm that enhances large language models by
    retrieving relevant information from external knowledge bases at
    inference time. It improves accuracy and credibility, especially for
    knowledge-intensive tasks, and allows knowledge to be updated without
    retraining.}
  \defn{zhao-2023}{A form of \emph{retrieval augmentation} for large language
    models that injects relevant external knowledge into the prompt to
    improve factuality and to give access to up-to-date or private
    information.}
  \note{RAG reduces hallucination but introduces new risks. Retrieved
    content can carry an indirect prompt injection.}
  \related{hallucination,prompt-injection,embedding}
\end{entry}
